\chapter{Análisis Tensorial}

\section{Convención de suma de Einstein}

Sea el espacio vectorial $V(n,K)$, generalización a $n$ dimensiones sobre
un campo $K$ del espacio vectorial estudiado en el capítulo anterior, con
base $\{\vec e_i\}_{i=1}^n$. Si $\vec v\in V(n,K)$, se escribe como
combinación lineal de la base elegida,
\[
\vec v = \sum_{i=1}^{n} v^i \vec e_i,
\]
donde $v^i$ es la componente del vector respecto de la base.

\begin{definicion}[Convención de Einstein]
Debido a que al trabajar con muchos índices aparecen muchas sumas,
adoptamos la \textbf{convención de suma de Einstein}: eliminaremos el
símbolo de suma $\sum$ y asumiremos que \textbf{índices repetidos
significan suma}, a menos que se diga explícitamente lo contrario.
\end{definicion}

Con esta convención, la ecuación anterior se simplifica a
\[
\vec v = v^i \vec e_i .
\]
Del mismo modo, si $\vec f \in V^*(n,K)$ es un vector dual escrito en la
base dual $\{\vec E^i\}_{i=1}^n$,
\[
\vec f = f_i \vec E^i .
\]
Los vectores del espacio dual son mapeos lineales que actúan sobre los
vectores del espacio vectorial y los llevan a $K$,
\[
f(v) = f_i v^i \in K,
\]
donde recordamos que los vectores duales se definen mediante
$E^i(e_j)=\delta^i_j$. A partir de aquí, salvo que se indique lo
contrario, se usará siempre la convención de Einstein.

\section{Subir y bajar índices}

Como se estudió en el capítulo anterior, la métrica $g_{ij}$ permite
definir el producto escalar y, además, permite pasar del espacio
vectorial al espacio vectorial dual,
\[
g_{ij} : V(n,K) \longrightarrow V^*(n,K), \qquad v_i = g_{ij} v^j .
\]
Como $\det(g_{ij})\neq 0$, existe la métrica inversa $g^{ij}$, tal que
$g^{ik}g_{kj}=\delta^i_j$, la cual permite pasar del espacio vectorial
dual al espacio vectorial,
\[
g^{ij} : V^*(n,K) \longrightarrow V(n,K), \qquad f^i = g^{ij} f_j .
\]

\begin{observacion}
Decimos que la métrica \textbf{baja índices}: convierte índices
contravariantes en covariantes. La métrica inversa \textbf{sube
índices}: convierte índices covariantes en contravariantes.
\end{observacion}

\begin{proposicionc}
La cantidad $v_i = g_{ij}v^j$, construida a partir del tensor $v^i$ de
orden $(1,0)$, es un tensor de orden $(0,1)$.
\end{proposicionc}

\begin{proof}
Sea la transformación de coordenadas $x'^i=x'^i(x^j)$. Usando la
definición $v_i=g_{ij}v^j$ en el sistema $O'x'^i$ y las leyes de
transformación de $g_{ij}$ y $v^j$,
\[
v'_i = g'_{ij} v'^j = \frac{\partial x^k}{\partial x'^i}\frac{\partial
x^l}{\partial x'^j} g_{kl}\, \frac{\partial x'^j}{\partial x^n} v^n =
\frac{\partial x^k}{\partial x'^i}\left(\frac{\partial x^l}{\partial
x'^j}\frac{\partial x'^j}{\partial x^n}\right) g_{kl} v^n = \frac{\partial
x^k}{\partial x'^i}\, \delta^l_n\, g_{kl} v^n = \frac{\partial
x^k}{\partial x'^i}\, g_{kl} v^l = \frac{\partial x^k}{\partial x'^i}
v_k,
\]
donde en el último paso se usó nuevamente $v_k=g_{kl}v^l$. Por lo tanto
\[
v'_i = \frac{\partial x^k}{\partial x'^i}\, v_k ,
\]
que es precisamente la ley de transformación de un tensor de orden
$(0,1)$. \qedhere
\end{proof}

\begin{proposicionc}
La cantidad $f^i=g^{ij}f_j$, construida a partir del tensor $f_i$ de
orden $(0,1)$, es un tensor de orden $(1,0)$.
\end{proposicionc}

\begin{proof}
Análogamente,
\[
f'^i = g'^{ij}f'_j = \frac{\partial x'^i}{\partial x^k}\frac{\partial
x'^j}{\partial x^l} g^{kl}\, \frac{\partial x^n}{\partial x'^j} f_n =
\frac{\partial x'^i}{\partial x^k}\left(\frac{\partial x'^j}{\partial
x^l}\frac{\partial x^n}{\partial x'^j}\right) g^{kl} f_n = \frac{\partial
x'^i}{\partial x^k}\, \delta^n_l\, g^{kl} f_n = \frac{\partial
x'^i}{\partial x^k}\, g^{kl}f_l = \frac{\partial x'^i}{\partial x^k}
f^k , \qedhere
\]
que es la ley de transformación de un tensor de orden $(1,0)$.
\end{proof}

\section{El producto tensorial}

\subsection{Motivación: producto tensorial de matrices columna}

Comencemos con $n=3$ y el espacio vectorial $\mathbb{R}^3$, escribiendo
los vectores como matrices columna.

\begin{definicion}[Producto tensorial de matrices columna]
Sean los vectores $v=(v^1,v^2,v^3)^T$ y $w=(w^1,w^2,w^3)^T$. Se define el
\textbf{producto tensorial} entre $v$ y $w$, denotado $\otimes$, como la
matriz
\[
M = v\otimes w = \begin{pmatrix} v^1 \\ v^2 \\ v^3 \end{pmatrix} \otimes
\begin{pmatrix} w^1 \\ w^2 \\ w^3 \end{pmatrix} = \begin{pmatrix}
v^1w^1 & v^1w^2 & v^1w^3 \\ v^2w^1 & v^2w^2 & v^2w^3 \\ v^3w^1 & v^3w^2 &
v^3w^3 \end{pmatrix}.
\]
\end{definicion}

\begin{ejemplo}
Sean $v=(1,-2,3)^T$ y $w=(1,1,2)^T$; entonces
\[
M = v\otimes w = \begin{pmatrix} 1 & 1 & 2 \\ -2 & -2 & -4 \\ 3 & 3 & 6
\end{pmatrix}.
\]
\end{ejemplo}

El producto tensorial entre dos matrices de $3\times 1$ genera una matriz
de $3\times 3$; como las matrices de $3\times 3$ forman un espacio
vectorial de dimensión $3\times 3=9$, su base canónica debe contener $9$
matrices linealmente independientes, dadas por
\[
e_{ij} = \left(\text{matriz con un } 1 \text{ en la fila } i,\ \text{columna } j,\ \text{y } 0
\text{ en el resto}\right), \qquad i,j=1,2,3.
\]
Cualquier matriz $M$ de $3\times 3$ se escribe entonces como combinación
lineal de estas 9 matrices base. Reescribiendo el producto tensorial
(11) en esta base:
\[
M = v\otimes w = v^1w^1e_{11} + v^1w^2e_{12} + \dots + v^3w^3e_{33} =
v^iw^j e_{ij} \qquad \text{(convención de Einstein)}.
\]
Por otro lado, $M$ se escribe en términos de la misma base como
$M = M^{ij}e_{ij}$, donde $M^{ij}$ es la componente de $M$ en la base
$\{e_{ij}\}$. Comparando ambas expresiones,
\[
M^{ij} = v^iw^j .
\]

\begin{observacion}
No todas las matrices de $3\times 3$ se pueden escribir en la forma
$M^{ij}=v^iw^j$: por ejemplo, la matriz
$\begin{pmatrix} 0 & -1 & 0 \\ 1 & 0 & 0 \\ 0 & 0 & 1\end{pmatrix}$
no admite vectores $v,w$ tales que $M=v\otimes w$. Sin embargo,
\textbf{todas} las matrices de $3\times3$, deriven o no de un producto
tensorial, se pueden escribir en la base $\{e_{ij}\}$; esta forma general
es la que se generaliza a continuación.
\end{observacion}

\subsection{Producto tensorial de espacios vectoriales genéricos}

El producto tensorial es un producto lineal:
\[
v\otimes(w+u) = v\otimes w + v\otimes u, \qquad (w+u)\otimes v = w\otimes
v + u\otimes v .
\]
Para $v,w\in V(n,K)$, el producto tensorial se escribe entonces
\[
M = v\otimes w = v^i\vec e_i \otimes w^j\vec e_j = v^iw^j\, \vec
e_i\otimes\vec e_j = M^{ij}\, \vec e_i\otimes\vec e_j, \qquad M^{ij}=
v^iw^j,
\]
y, comparando con el caso de matrices columna, se concluye que
\[
e_{ij} = \vec e_i \otimes \vec e_j, \qquad i,j=1,\dots,n.
\]

\begin{definicion}[Componente irreducible]
La componente $M^{ij}$ de la matriz $M$ se dice \textbf{irreducible} si
\emph{no} se puede escribir en la forma $M^{ij}=v^iw^j$, es decir, si
$M \neq v\otimes w$.
\end{definicion}

\begin{definicion}[Espacio $V(n,K)\otimes V(n,K)$]
Se define el espacio vectorial $V(n,K)\otimes V(n,K)$, con base
$\{\vec e_i\otimes\vec e_j\}_{i,j=1}^n$, tal que un vector $T$ de este
espacio se escribe como
\[
T = T^{ij}\, \vec e_i \otimes \vec e_j ,
\]
con $T^{ij}$ la componente de $T$ en la base elegida.
\end{definicion}

Bajo una transformación de coordenadas $x'^i=x'^i(x^j)$, la base
$\{\vec e_i\}$ transforma como $\vec e\,'_i = (\partial x^j/\partial
x'^i)\, \vec e_j$. Un vector $T\in V(n,K)\otimes V(n,K)$ escrito en los
sistemas $Ox^i$ y $O'x'^i$ debe satisfacer el \textbf{principio de
covariancia}, $T\big|_{O'}=T\big|_O$:

\begin{proposicionc}
Bajo la transformación de coordenadas $x'^i=x'^i(x^j)$, la componente de
un vector $T\in V(n,K)\otimes V(n,K)$ transforma como
\[
T'^{ij} = \frac{\partial x'^i}{\partial x^k}\frac{\partial
x'^j}{\partial x^l}\, T^{kl} .
\]
Es decir, $T^{ij}$ es un tensor de orden $(2,0)$.
\end{proposicionc}

\begin{proof}
De $T^{ij}\vec e_i\otimes\vec e_j = T'^{kl}\vec e\,'_k\otimes\vec e\,'_l
= T'^{kl}\dfrac{\partial x^i}{\partial x'^k}\dfrac{\partial
x^j}{\partial x'^l}\, \vec e_i\otimes\vec e_j$ y la independencia lineal
de la base $\{\vec e_i\otimes\vec e_j\}$, se obtiene
\[
T^{ij} = \frac{\partial x^i}{\partial x'^k}\frac{\partial
x^j}{\partial x'^l}\, T'^{kl},
\]
de donde, despejando $T'^{ij}$ (transformación inversa), se sigue el
resultado. \qedhere
\end{proof}

\subsection{Generalización a tensores de orden \texorpdfstring{$(p,0)$}
{(p,0)}}

\begin{definicion}
Se define el espacio vectorial $\underbrace{V(n,K)\otimes\dots\otimes
V(n,K)}_{p\text{ veces}}$, con base $\{\vec e_{i_1}\otimes\dots\otimes
\vec e_{i_p}\}_{i_1,\dots,i_p=1}^n$, tal que un vector $T$ de este
espacio se escribe
\[
T = T^{i_1i_2\dots i_p}\, \vec e_{i_1}\otimes\vec e_{i_2}\otimes\dots
\otimes\vec e_{i_p}.
\]
\end{definicion}

\begin{proposicionc}
Bajo $x'^i=x'^i(x^j)$, la componente $T^{i_1\dots i_p}$ transforma como
\[
T'^{i_1i_2\dots i_p} = \frac{\partial x'^{i_1}}{\partial x^{j_1}}
\frac{\partial x'^{i_2}}{\partial x^{j_2}} \cdots \frac{\partial
x'^{i_p}}{\partial x^{j_p}}\, T^{j_1j_2\dots j_p}.
\]
\end{proposicionc}

\begin{definicion}[Tensor contravariante de orden $p$]
La cantidad $T^{i_1\dots i_p}$ se dice un \textbf{tensor contravariante
de orden $p$}, o tensor de orden $(p,0)$, si bajo una transformación de
coordenadas transforma como en la proposición anterior. Su ley de
transformación contiene $p$ matrices $\partial x'^i/\partial x^j$ de
forma homogénea (todas del mismo tipo).
\end{definicion}

\subsection{Generalización a tensores de orden \texorpdfstring{$(0,q)$}
{(0,q)}}

De manera análoga, usando el espacio vectorial dual $V^*(n,K)$: para
$V,W\in V^*(n,K)$,
\[
T = V\otimes W = V_iE^i \otimes W_jE^j = V_iW_j\, E^i\otimes E^j =:
T_{ij}\, E^i\otimes E^j, \qquad T_{ij}=V_iW_j.
\]
Nótese que $T\notin V^*(n,K)$, sino que pertenece a un nuevo espacio
vectorial dual, y en general $T_{ij}\neq V_iW_j$ para un tensor
arbitrario.

\begin{definicion}
Se define el espacio vectorial $\underbrace{V^*(n,K)\otimes\dots\otimes
V^*(n,K)}_{q\text{ veces}}$, con base $\{E^{i_1}\otimes\dots\otimes
E^{i_q}\}_{i_1,\dots,i_q=1}^n$, tal que $T$ en este espacio se escribe
\[
T = T_{i_1i_2\dots i_q}\, E^{i_1}\otimes E^{i_2}\otimes\dots\otimes
E^{i_q}.
\]
\end{definicion}

\begin{proposicionc}
Bajo $x'^i=x'^i(x^j)$, la componente $T_{i_1\dots i_q}$ transforma como
\[
T'_{i_1i_2\dots i_q} = \frac{\partial x^{j_1}}{\partial x'^{i_1}}
\frac{\partial x^{j_2}}{\partial x'^{i_2}} \cdots \frac{\partial
x^{j_q}}{\partial x'^{i_q}}\, T_{j_1j_2\dots j_q}.
\]
\end{proposicionc}

\begin{definicion}[Tensor covariante de orden $q$]
La cantidad $T_{i_1\dots i_q}$ se dice un \textbf{tensor covariante de
orden $q$}, o tensor de orden $(0,q)$, si transforma como en la
proposición anterior; su ley de transformación contiene $q$ matrices
$\partial x^j/\partial x'^i$ de forma homogénea.
\end{definicion}

Puesto que $V^*(n,K)\otimes\dots\otimes V^*(n,K)$ (con $q$ factores) es
un espacio vectorial dual, actúa sobre $V(n,K)\otimes\dots\otimes
V(n,K)$ ($p$ factores, con $p=q$, sin índices libres): si
$V\in V(n,K)^{\otimes q}$ y $T\in V^*(n,K)^{\otimes q}$,
\[
T(V) = T_{i_1\dots i_q}\, V^{i_1\dots i_q} \in K .
\]

\subsection{Generalización a tensores de orden \texorpdfstring{$(1,1)$}
{(1,1)} y \texorpdfstring{$(p,q)$}{(p,q)}}

Para $V\in V(n,K)$ y $W\in V^*(n,K)$,
\[
T = V\otimes W = V^i\vec e_i\otimes W_jE^j = V^iW_j\, \vec e_i\otimes E^j
=: T^i_{\ j}\, \vec e_i\otimes E^j, \qquad T^i_{\ j}=V^iW_j,
\]
donde $T$ pertenece a un nuevo \textbf{espacio vectorial mixto}
$V(n,K)\otimes V^*(n,K)$ (en general $T^i_{\ j}\neq V^iW_j$).

\begin{proposicionc}
Si $T\in V(n,K)\otimes V^*(n,K)$, entonces bajo $x'^i=x'^i(x^j)$ la
componente $T^i_{\ j}$ transforma como
\[
T'^i_{\ j} = \frac{\partial x'^i}{\partial x^k}\frac{\partial
x^l}{\partial x'^j}\, T^k_{\ l} .
\]
\end{proposicionc}

\begin{definicion}[Tensor de orden $(1,1)$]
La cantidad $T^i_{\ j}$ se dice un tensor de orden $(1,1)$
(contravariante de orden uno y covariante de orden uno) si transforma
como en la proposición anterior.
\end{definicion}

\begin{observacion}
El espacio mixto se escribe convencionalmente $V(n,K)\otimes V^*(n,K)$ y
no al revés: como $V^*(n,K)$ es un espacio de funcionales lineales que
actúan sobre $V(n,K)$, escribir $V^*(n,K)\otimes V(n,K)$ podría sugerir
erróneamente que el dual actúa sobre el factor $V(n,K)$. (Alguna
literatura usa el orden opuesto; no debe generar confusión.) En la ley de
transformación de un tensor $(1,1)$ aparecen ambos tipos de matrices,
$\partial x'^i/\partial x^j$ y $\partial x^j/\partial x'^i$, pero
\emph{no} contraídas entre sí, pues actúan sobre índices distintos.
\end{observacion}

La generalización a un tensor de orden $(p,q)$ arbitrario es directa:

\begin{definicion}[Tensor de orden $(p,q)$]
Se define el espacio vectorial
\[
\underbrace{V(n,K)\otimes\dots\otimes V(n,K)}_{p\text{ veces}} \otimes
\underbrace{V^*(n,K)\otimes\dots\otimes V^*(n,K)}_{q\text{ veces}},
\]
con base $\{\vec e_{i_1}\otimes\dots\otimes\vec e_{i_p}\otimes
E^{j_1}\otimes\dots\otimes E^{j_q}\}$, tal que $T$ en este espacio se
escribe
\[
T = T^{i_1\dots i_p}_{\quad\ j_1\dots j_q}\, \vec e_{i_1}\otimes\dots
\otimes\vec e_{i_p}\otimes E^{j_1}\otimes\dots\otimes E^{j_q}.
\]
La cantidad $T^{i_1\dots i_p}_{\quad\ j_1\dots j_q}$ se dice un
\textbf{tensor de orden $(p,q)$} (contravariante de orden $p$ y
covariante de orden $q$) si bajo $x'^i=x'^i(x^j)$ transforma como
\[
T'^{i_1\dots i_p}_{\quad\ j_1\dots j_q} = \frac{\partial
x'^{i_1}}{\partial x^{k_1}}\cdots\frac{\partial x'^{i_p}}{\partial
x^{k_p}}\, \frac{\partial x^{l_1}}{\partial x'^{j_1}}\cdots\frac{\partial
x^{l_q}}{\partial x'^{j_q}}\, T^{k_1\dots k_p}_{\quad\ l_1\dots l_q},
\]
con $p$ matrices $\partial x'^i/\partial x^j$ y $q$ matrices $\partial
x^j/\partial x'^i$.
\end{definicion}

\begin{observacion}[Definición cualitativa]
Un tensor es una entidad algebraica de varias componentes que generaliza
los conceptos de escalar, vector y matriz de forma independiente del
sistema de coordenadas elegido: es \textbf{invariante en forma} ante un
cambio de coordenadas, $T\big|_O = T\big|_{O'}$.
\end{observacion}

\section{Subir y bajar índices en tensores generales}

La métrica, al actuar sobre un tensor de orden $(0,p)$ (o, en general,
sobre cualquier índice contravariante de un tensor de orden $(p,q)$),
genera la cantidad
\[
T^{i_1\dots i_{s-1}}_{\quad\quad\ \ i_{s+1}\dots i_p\, j} = g_{jk}\,
T^{i_1\dots i_{s-1}k i_{s+1}\dots i_p},
\]
que es un tensor de orden $(p-1,1)$ (se ``bajó'' un índice). En general
se pueden bajar $r$ índices con $r$ métricas:
\[
T^{i_1\dots i_s}_{\quad\ \ i_{s+r+1}\dots i_p\, k_1\dots k_r} =
g_{k_1j_1}\cdots g_{k_rj_r}\, T^{i_1\dots i_s j_1\dots j_r
i_{s+r+1}\dots i_p},
\]
que es un tensor de orden $(p-r,r)$. De modo análogo, la métrica inversa
permite \textbf{subir} $r$ índices covariantes de un tensor de orden
$(0,q)$, obteniendo un tensor de orden $(r,q-r)$.

\subsection{Los símbolos de Christoffel}

\begin{definicion}[Símbolos de Christoffel]
Los \textbf{símbolos de Christoffel} $\Gamma^k_{ij}$ se definen a
través de la derivada de la base coordenada,
\[
\frac{\partial \vec e_i}{\partial x^j} = \Gamma^k_{ij}\, \vec e_k .
\]
Miden el cambio de dirección (y de magnitud) de los vectores base al ser
trasladados a lo largo de curvas del espacio.
\end{definicion}

\begin{proposicionc}
Bajo la transformación de coordenadas $x'^i=x'^i(x^j)$, los símbolos de
Christoffel transforman como
\[
\Gamma'^{\,k}_{ij} = \frac{\partial x'^k}{\partial x^m}\frac{\partial
x^l}{\partial x'^i}\frac{\partial x^n}{\partial x'^j}\, \Gamma^m_{ln} +
\frac{\partial x'^k}{\partial x^m}\, \frac{\partial^2 x^m}{\partial
x'^i\partial x'^j} .
\]
\end{proposicionc}

\begin{observacion}
Los símbolos de Christoffel \textbf{no son un tensor}: su ley de
transformación contiene el término inhomogéneo
$(\partial x'^k/\partial x^m)(\partial^2 x^m/\partial x'^i\partial
x'^j)$, que depende del sistema de coordenadas, mientras que los
tensores, al ser covariantes, no dependen de él. Si los símbolos de
Christoffel fueran un tensor, transformarían únicamente con el primer
término (homogéneo),
$\Gamma'^{\,k}_{ij}=(\partial x'^k/\partial x^m)(\partial x^l/\partial
x'^i)(\partial x^n/\partial x'^j)\Gamma^m_{ln}$.

El primer término (homogéneo) mide el cambio \emph{covariante} de los
vectores base al ser trasladados, lo cual no depende del sistema
coordenado sino del tipo de espacio (curvo o plano) en el que se
trasladan. El segundo término (inhomogéneo) mide el cambio de los
vectores base al ser trasladados sobre líneas coordenadas curvas, y sí
depende del sistema coordenado elegido.
\end{observacion}

\section{Operaciones con tensores}

Los tensores viven en espacios vectoriales, por lo que satisfacen las
operaciones propias de un espacio vectorial.

\begin{definicion}[Suma]
Sean $V^{i_1\dots i_p}_{\ \ j_1\dots j_q}$ y $W^{i_1\dots
i_p}_{\ \ j_1\dots j_q}$ dos tensores del mismo orden $(p,q)$ (por
homogeneidad, deben serlo). Se define la suma
\[
T^{i_1\dots i_p}_{\quad\ j_1\dots j_q} = \alpha\, V^{i_1\dots
i_p}_{\quad\ j_1\dots j_q} + \beta\, W^{i_1\dots i_p}_{\quad\ j_1\dots
j_q}, \qquad \alpha,\beta\in K.
\]
\end{definicion}

\begin{definicion}[Producto tensorial]
El producto tensorial entre un tensor $V^{i_1\dots
i_p}_{\quad\ j_1\dots j_q}$ de orden $(p,q)$ y un tensor $W^{i_1\dots
i_{p'}}_{\quad\ j_1\dots j_{q'}}$ de orden $(p',q')$ se define como
\[
M^{i_1\dots i_{p+p'}}_{\quad\quad\ \ j_1\dots j_{q+q'}} = V^{i_1\dots
i_p}_{\quad\ j_1\dots j_q}\, W^{i_{p+1}\dots i_{p+p'}}_{\quad\quad\
j_{q+1}\dots j_{q+q'}} .
\]
\end{definicion}

\begin{proposicionc}
El producto tensorial de un tensor de orden $(p,q)$ con uno de orden
$(p',q')$ es un tensor de orden $(p+p',q+q')$.
\end{proposicionc}

Por lo tanto, el producto tensorial permite \textbf{construir tensores de
mayor orden}. La operación complementaria, llamada \textbf{contracción},
permite construir tensores de \textbf{menor orden}.

\begin{ejemplo}
Sea el tensor $K^i_{\ jl}$. Sumando (contrayendo) los índices $i$ y $l$,
\[
K^l_{\ jl} \qquad \text{(convención de Einstein: suma sobre } l\text{)},
\]
se obtiene una cantidad con un solo índice libre, $j$.
\end{ejemplo}

\begin{proposicionc}
Sea el tensor $K^i_{\ jl}$ de orden $(1,2)$; la contracción $K^l_{\ jl}$
es un tensor de orden $(0,1)$.
\end{proposicionc}

\begin{definicion}[Contracción]
La \textbf{contracción} consiste en poner bajo la operación de suma la
misma cantidad de índices contravariantes y covariantes.
\end{definicion}

\begin{proposicionc}
Sea $V^{i_1\dots i_p}_{\quad\ j_1\dots j_q}$ un tensor de orden $(p,q)$.
Efectuando una contracción de $r$ índices contravariantes con $r$
índices covariantes ($r<\min\{p,q\}$),
\[
V^{i_1\dots i_s k_1\dots k_r\, i_{s+r+1}\dots i_p}_{\quad\quad\quad\quad\
\ \, j_1\dots j_s k_1\dots k_r\, j_{s+r+1}\dots j_q},
\]
se obtiene un tensor de orden $(p-r,q-r)$.
\end{proposicionc}

\begin{observacion}
Si se desea contraer dos índices de la misma naturaleza (por ejemplo, dos
índices covariantes $j,k$ del tensor $T_{ijk}$), esto no cumple la
definición de contracción directamente. La solución es subir uno de los
dos índices con la métrica inversa antes de contraer:
\[
T_i^{\ k} = g^{kl}T_{ijl}\, \delta^j_{(\cdot)} \quad\longrightarrow\quad
\text{primero } T^{\ \ k}_{ij} = g^{kl}T_{ijl}, \quad\text{luego se
contrae } i \text{ con } k.
\]
\end{observacion}

\section{Simetrías}

\begin{definicion}[Tensor simétrico]
Un tensor se dice \textbf{simétrico} si bajo el intercambio de dos
índices el tensor no cambia:
\[
T^{\dots i\dots j\dots} = +T^{\dots j\dots i\dots}, \qquad T_{\dots
i\dots j\dots} = +T_{\dots j\dots i\dots}.
\]
\end{definicion}

\begin{ejemplo}
La métrica y la métrica inversa son simétricas: $g_{ij}=g_{ji}$,
$g^{ij}=g^{ji}$.
\end{ejemplo}

\begin{observacion}
La simetría de un tensor afecta el número de componentes independientes.
Por ejemplo, la métrica $g_{ij}$ posee, en principio, $n\times n=n^2$
componentes, pero al ser simétrica solo tiene $n(n+1)/2$ componentes
independientes.
\end{observacion}

\begin{proposicionc}
Un tensor de orden $(p,0)$, $T^{i_1\dots i_p}$, posee en principio $n^p$
componentes independientes. Si dos de sus índices son simétricos,
\[
T^{i_1\dots j\dots k\dots i_p} = +T^{i_1\dots k\dots j\dots i_p},
\]
el número de componentes independientes se reduce a
\[
n^{p-2}\times \frac{n(n+1)}{2} = \frac{1}{2}n^{p-1}(n+1).
\]
\end{proposicionc}

La misma conclusión se obtiene para tensores de orden $(0,q)$.

\begin{definicion}[Tensor antisimétrico]
Un tensor se dice \textbf{antisimétrico} si bajo el intercambio de dos
índices el tensor cambia de signo:
\[
T^{\dots i\dots j\dots} = -T^{\dots j\dots i\dots}, \qquad T_{\dots
i\dots j\dots} = -T_{\dots j\dots i\dots}.
\]
\end{definicion}

\begin{proposicionc}
Un tensor de orden $(2,0)$ antisimétrico, $T^{ij}=-T^{ji}$, en un espacio
plano de $n$ dimensiones posee $n(n-1)/2$ componentes independientes (en
vez de $n^2$). En general, para un tensor de orden $(p,0)$ con dos
índices antisimétricos,
\[
n^{p-2}\times\frac{n(n-1)}{2} = \frac{1}{2}n^{p-1}(n-1).
\]
La misma conclusión se obtiene para tensores de orden $(0,q)$.
\end{proposicionc}

\begin{proposicionc}
Un tensor antisimétrico posee ceros en los índices que se repiten:
$T^{\dots i\dots i\dots}=0$ (sin suma).
\end{proposicionc}

\begin{proof}
Haciendo $i=j$ en la definición de antisimetría (sin sumar sobre $i$),
$T^{\dots i\dots i\dots} = -T^{\dots i\dots i\dots}$, de donde
$T^{\dots i\dots i\dots}=0$. \qedhere
\end{proof}

\begin{definicion}[Descomposición simétrica-antisimétrica]
Dado un tensor $T^{ij}$, siempre es posible construir su parte simétrica
y antisimétrica:
\[
T^{(ij)} := \frac{1}{2}\left(T^{ij}+T^{ji}\right) = T^{(ji)}, \qquad
T^{[ij]} := \frac{1}{2}\left(T^{ij}-T^{ji}\right) = -T^{[ji]} .
\]
El factor $1/2$ evita contar dos veces. Lo mismo se aplica a $T_{ij}$.
\end{definicion}

\begin{proposicionc}
Todo tensor $T^{ij}$ (o $T_{ij}$) se descompone como
\[
T^{ij} = T^{(ij)} + T^{[ij]} .
\]
\end{proposicionc}

\begin{proposicionc}
Si $A^{ij}$ es simétrico y $B_{ij}$ es antisimétrico, entonces
$A^{ij}B_{ij}=0$.
\end{proposicionc}

\begin{proposicionc}
Si $A^{ij}$ es antisimétrico y $B_{ij}$ es un tensor cualquiera, entonces
$A^{ij}B_{ij} = A^{ij}B_{[ij]}$.
\end{proposicionc}

\begin{ejemplo}
Sea $T^{jkl}_{\ \ \ i}$; formemos sus partes simétrica y antisimétrica en
los índices $i,k$. Como $i$ es covariante y $k$ contravariante, son de
distinta naturaleza y deben igualarse antes de (anti)simetrizar. Subiendo
$i$ con la métrica inversa, $T^{ijkl}=g^{in}T^{jkl}_{\ \ \ n}$:
\begin{align*}
T^{(i|j|k)l} &= \frac{1}{2}\left(T^{ijkl}+T^{kjil}\right) =
\frac{1}{2}\left(g^{in}T^{jkl}_{\ \ \ n} + g^{kn}T^{jil}_{\ \ \ n}\right)
= g^{n(i|}T^{j|k)l}_{\quad\ n}, \\
T^{[i|j|k]l} &= \frac{1}{2}\left(T^{ijkl}-T^{kjil}\right) = g^{n[i|}
T^{j|k]l}_{\quad\ n},
\end{align*}
donde las barras verticales indican que el índice $j$ queda excluido de
la (anti)simetrización. Alternativamente, bajando el índice $k$ con la
métrica $g_{ij}$ se obtiene la descomposición equivalente
$T^{j\ \ l}_{\ (i\,k)}$, $T^{j\ \ l}_{\ [i\,k]}$.
\end{ejemplo}

\begin{definicion}[Tensor completamente antisimétrico]
Un tensor se dice \textbf{completamente antisimétrico} si cambia de
signo bajo cualquier permutación de un par de índices.
\end{definicion}

El caso completamente antisimétrico es de particular interés (más que el
completamente simétrico). Para un tensor de orden $(0,q)$ completamente
antisimétrico, $T_{i_1i_2\dots i_q}$, el número de componentes
independientes es
\[
\binom{n}{q} = \frac{n!}{q!(n-q)!} .
\]

\begin{observacion}
Si $q>n$ (más índices que dimensiones), entonces $\binom{n}{q}=0$: el
tensor es idénticamente nulo, pues los índices deben repetirse
obligatoriamente y, como vimos, los índices repetidos anulan un tensor
completamente antisimétrico.
\end{observacion}

Si $q=n$, el tensor completamente antisimétrico posee $\binom{n}{n}=1$
componente independiente.

\subsection{El símbolo de Levi-Civita}

\begin{ejemplo}[Caso $n=2$]
El tensor $T_{ij}$ completamente antisimétrico en $2$ dimensiones,
$T_{ij}=-T_{ji}$, se escribe
\[
T_{ij} = \begin{pmatrix} 0 & T_{12} \\ -T_{12} & 0 \end{pmatrix} =
T_{12} \begin{pmatrix} 0 & 1 \\ -1 & 0 \end{pmatrix} = T_{12}\,
\epsilon_{ij},
\]
posee una sola componente independiente $T_{12}$. Se define
\[
\epsilon_{ij} = \begin{cases} +1, & (ij)=(12) \\ -1, & (ij)=(21) \\ 0, &
i=j \end{cases}
\]
llamado \textbf{símbolo de Levi-Civita} en dos dimensiones.
\end{ejemplo}

\begin{observacion}
El símbolo de Levi-Civita \textbf{no} es un tensor de orden $(0,2)$, sino
una \emph{densidad tensorial} de orden $(0,2)$ y peso $-1$.
\end{observacion}

\begin{ejemplo}[Caso $n=3$]
El tensor totalmente antisimétrico $T_{ijk}$ en tres dimensiones posee
una única componente independiente $T_{123}$, tal que
\[
T_{ijk} = T_{123}\, \epsilon_{ijk},
\]
donde $\epsilon_{ijk}$ es el símbolo de Levi-Civita en tres dimensiones,
\[
\epsilon_{ijk} = \begin{cases} +1, & (ijk) \text{ permutación par de }
(123) \\ -1, & (ijk) \text{ permutación impar de } (123) \\ 0, & \text{en
otro caso.} \end{cases}
\]
Por ejemplo, $\epsilon_{213}=-\epsilon_{123}=-1$,
$\epsilon_{132}=-\epsilon_{123}=-1$, $\epsilon_{312} =
-\epsilon_{132}=+1$. Los signos, usando índices simbólicos $ijk$ en vez
de $123$, se obtienen por \textbf{permutación cíclica}: por ejemplo,
$\epsilon_{jik}=-\epsilon_{ijk}$, $\epsilon_{ikj}=-\epsilon_{ijk}$,
$\epsilon_{kji} = -\epsilon_{jki} = \epsilon_{jik}=-\epsilon_{ijk}$.
\end{ejemplo}

Las propiedades siguientes se establecen en un sistema coordenado
cartesiano de un espacio plano, donde $g_{ij}=\delta_{ij}$.

\begin{proposicionc}
El símbolo de Levi-Civita $\epsilon_{ijk}$ se puede escribir como
\[
\epsilon_{ijk} = \det\begin{pmatrix} \delta_{i1} & \delta_{i2} &
\delta_{i3} \\ \delta_{j1} & \delta_{j2} & \delta_{j3} \\ \delta_{k1} &
\delta_{k2} & \delta_{k3} \end{pmatrix},
\]
y satisface
\[
\epsilon_{ijk}\epsilon_{lmn} = \det\begin{pmatrix} \delta_{il} &
\delta_{im} & \delta_{in} \\ \delta_{jl} & \delta_{jm} & \delta_{jn} \\
\delta_{kl} & \delta_{km} & \delta_{kn} \end{pmatrix}, \qquad
\epsilon_{ijk}\epsilon_{lmk} = \delta_{il}\delta_{jm} -
\delta_{im}\delta_{jl}, \qquad \epsilon_{ijk}\epsilon_{ljk} =
2\delta_{il} .
\]
\end{proposicionc}

\begin{proposicionc}
El determinante de un tensor $A^{ij}$ se puede escribir como
\[
\det(A^{ij}) = \epsilon_{ijk}A^{i1}A^{j2}A^{k3} = \frac{1}{3!}\,
\epsilon_{ijk}\epsilon_{lmn}\, A^{il}A^{jm}A^{kn} .
\]
\end{proposicionc}

\section{Producto vectorial}

La base canónica $\{\ihat,\jhat,\khat\}$ de $\mathbb{R}^3$ satisface
$\ihat\times\jhat=\khat$, $\ihat\times\khat=-\jhat$,
$\jhat\times\khat=\ihat$. Escribiendo $\vec e_1=\ihat$, $\vec
e_2=\jhat$, $\vec e_3=\khat$, esta álgebra se resume como
\[
\vec e_i \times \vec e_j = \epsilon_{ijk}\, \vec e_k .
\]

\begin{proposicionc}
El producto cruz entre $\vec A$ y $\vec B$, $\vec C = \vec A\times\vec
B$, se escribe en notación de índices como
\[
C^i = \epsilon^{ijk}A_jB_k .
\]
\end{proposicionc}

\begin{proof}
\[
\vec C = \vec A\times\vec B = A^i\vec e_i\times B^j\vec e_j = A^iB^j\,
\vec e_i\times\vec e_j = A^iB^j\epsilon_{ijk}\vec e_k = \epsilon_{jki}
A^jB^k \vec e_i = \epsilon_{ijk}A^jB^k\, \vec e_i,
\]
usando la simetría cíclica $\epsilon_{jki}=\epsilon_{ijk}$. Comparando
con $\vec C = C^i\vec e_i$, se concluye $C^i = \epsilon^{ijk}A_jB_k$
(en cartesianas, $\epsilon^{ijk}=\epsilon_{ijk}$ e índices arriba/abajo
son intercambiables). \qedhere
\end{proof}

\begin{ejemplo}
Desarrollando explícitamente para $i=1,2,3$ (usando
$\epsilon_{1jk},\epsilon_{2jk},\epsilon_{3jk}$ y anulando los términos
con índices repetidos) se recupera la fórmula usual del producto cruz:
\[
C^1 = A^2B^3-A^3B^2, \qquad C^2 = A^3B^1-A^1B^3, \qquad C^3 =
A^1B^2-A^2B^1.
\]
\end{ejemplo}

\begin{proposicionc}
El vector $\vec C=\vec A\times\vec B$ es perpendicular a $\vec A$ y a
$\vec B$.
\end{proposicionc}

\begin{proof}
Calculemos $\vec C\cdot\vec A = C^iA_i = \epsilon^{ijk}A_jB_kA_i =
\epsilon^{ijk}A_iA_jB_k$. Definiendo el tensor simétrico
$T^{ij}=A^iA^j=A^jA^i=T^{ji}$,
\[
\vec C\cdot\vec A = \epsilon^{ijk}T^{ij}B_k .
\]
Por ser $T^{ij}$ simétrico y $\epsilon^{ijk}$ antisimétrico en $(i,j)$,
esta contracción se anula (Proposición sobre $A^{ij}B_{ij}=0$ para $A$
simétrico y $B$ antisimétrico):
\[
\epsilon^{ijk}T^{ij}B_k = -\epsilon^{jik}T^{ij}B_k = -\epsilon^{ijk}
T^{ji}B_k = -\epsilon^{ijk}T^{ij}B_k \;\Rightarrow\; \vec C\cdot\vec A =
0 .
\]
La misma conclusión se obtiene para $\vec C\cdot\vec B=0$. \qedhere
\end{proof}

Se sigue que $\vec A\times\vec B$ es perpendicular al plano que
contienen $\vec A$ y $\vec B$, y por lo tanto todo vector contenido en
dicho plano es perpendicular a $\vec A\times\vec B$: en efecto, todo
vector $\vec D$ del plano generado por $\vec A$ y $\vec B$ se escribe
\[
\vec D = \frac{D}{\sin(\theta_1+\theta_2)}\left(\sin\theta_2\, \hat A +
\sin\theta_1\, \hat B\right),
\]
con $D=\|\vec D\|$, $\hat A=\vec A/A$, $\hat B=\vec B/B$, de donde
$\vec D\cdot\vec C = 0$ directamente.

\begin{proposicionc}
Si $\vec A,\vec B,\vec C$ están contenidos en un plano, entonces
$\epsilon_{ijk}A^iB^jC^k=0$. Si son no coplanares, el volumen que forman
es
\[
V = \epsilon_{ijk}A^iB^jC^k .
\]
\end{proposicionc}

\begin{proposicionc}
La norma de $\vec C=\vec A\times\vec B$ es
\[
\|\vec C\| = \|\vec A\|\,\|\vec B\|\,\sin\theta ,
\]
donde $\theta$ es el ángulo entre $\vec A$ y $\vec B$.
\end{proposicionc}

\begin{proof}
\begin{align*}
\|\vec C\|^2 &= C^iC_i = \epsilon^{ijk}A_jB_k\,\epsilon_{inm}A^nB^m \\
&= \left(\delta^j_n\delta^k_m - \delta^j_m\delta^k_n\right) A_jB_kA^nB^m
\\
&= A_jA^jB_kB^k - A_jB^jA^kB_k \\
&= \|\vec A\|^2\|\vec B\|^2 - (\vec A\cdot\vec B)^2,
\end{align*}
usando la identidad $\epsilon^{ijk}\epsilon_{inm} =
\delta^j_n\delta^k_m-\delta^j_m\delta^k_n$. Como $\vec A\cdot\vec B =
\|\vec A\|\|\vec B\|\cos\theta$,
\[
\|\vec C\|^2 = \|\vec A\|^2\|\vec B\|^2(1-\cos^2\theta) = \|\vec
A\|^2\|\vec B\|^2\sin^2\theta \;\Rightarrow\; \|\vec C\| = \|\vec
A\|\,\|\vec B\|\,\sin\theta . \qedhere
\]
\end{proof}

\section{Derivada covariante}

Sea $\vec A$ un vector escrito en la base coordenada $Ox^i$. Calculemos
su derivada parcial $\partial_i := \partial/\partial x^i$:
\[
\partial_i \vec A = \partial_i(A^j\vec e_j) = (\partial_iA^j)\vec e_j +
A^j(\partial_i\vec e_j) = (\partial_iA^j)\vec e_j + A^j\Gamma^k_{ij}\vec
e_k = \left(\partial_iA^j + \Gamma^j_{ik}A^k\right)\vec e_j ,
\]
donde se usó la definición de los símbolos de Christoffel y se
renombraron índices mudos.

\begin{definicion}[Derivada covariante de un vector]
Sea $A^i$ un tensor de orden $(1,0)$. Se define su \textbf{derivada
covariante} como
\[
\nabla_i A^j := \partial_i A^j + \Gamma^j_{ik}A^k .
\]
\end{definicion}

Con esta definición, $\partial_i\vec A = \nabla_iA^j\,\vec e_j$: la
derivada covariante $\nabla_iA^j$ reemplaza, en notación tensorial, a la
derivada parcial $\partial_iA^j$.

\begin{proposicionc}
Si $A^i$ es un tensor de orden $(1,0)$, entonces $\nabla_iA^j$ es un
tensor de orden $(1,1)$.
\end{proposicionc}

Análogamente, para un vector dual $\vec f\in V^*(n,K)$,
\[
\partial_i\vec f = \partial_i(f_jE^j) = (\partial_if_j)E^j +
f_j\partial_iE^j = (\partial_if_j)E^j - f_k\Gamma^k_{ij}E^j =
\left(\partial_if_j - \Gamma^k_{ij}f_k\right)E^j ,
\]
donde se usó $\partial_iE^j = -\Gamma^j_{ik}E^k$ (consecuencia de derivar
$E^i(e_j)=\delta^i_j$).

\begin{definicion}[Derivada covariante de un covector]
Sea $A_i$ un tensor de orden $(0,1)$. Se define
\[
\nabla_iA_j := \partial_iA_j - \Gamma^k_{ij}A_k .
\]
\end{definicion}

\begin{proposicionc}
Si $A_i$ es un tensor de orden $(0,1)$, entonces $\nabla_iA_j$ es un
tensor de orden $(0,2)$.
\end{proposicionc}

\begin{observacion}
Resumiendo: $\nabla_iA^j = \partial_iA^j+\Gamma^j_{ik}A^k$ (el índice
arriba aparece con $+\Gamma$); $\nabla_iA_j =
\partial_iA_j-\Gamma^k_{ij}A_k$ (el índice abajo aparece con $-\Gamma$).
\end{observacion}

Para un tensor mixto $A\in V(n,\mathbb{R})\otimes V^*(n,\mathbb{R})$, un
cálculo análogo (derivando $A=A^j_{\ k}\,\vec e_j\otimes E^k$ y usando
las leyes de derivación de ambas bases) da
\[
\partial_iA = \left(\partial_iA^j_{\ k} + \Gamma^j_{il}A^l_{\ k} -
\Gamma^l_{ik}A^j_{\ l}\right)\vec e_j\otimes E^k .
\]

\begin{definicion}[Derivada covariante de un tensor $(1,1)$]
Sea $A^i_{\ j}$ un tensor de orden $(1,1)$. Se define
\[
\nabla_iA^j_{\ k} = \partial_iA^j_{\ k} + \Gamma^j_{il}A^l_{\ k} -
\Gamma^l_{ik}A^j_{\ l} .
\]
\end{definicion}

\begin{proposicionc}
Si $A^i_{\ j}$ es un tensor de orden $(1,1)$, entonces $\nabla_iA^j_{\
k}$ es un tensor de orden $(1,2)$.
\end{proposicionc}

La generalización a tensores de orden $(p,q)$ arbitrario es directa:

\begin{definicion}[Derivada covariante general]
Sea $A^{i_1\dots i_p}_{\quad\ j_1\dots j_q}$ un tensor de orden $(p,q)$.
Se define su derivada covariante como
\[
\nabla_i A^{j_1\dots j_p}_{\quad\ k_1\dots k_q} = \partial_i A^{j_1\dots
j_p}_{\quad\ k_1\dots k_q} + \Gamma^{j_1}_{il}A^{lj_2\dots j_p}_{\quad\
\ k_1\dots k_q} + \dots + \Gamma^{j_p}_{il}A^{j_1\dots j_{p-1}l}_{\quad\
\quad\ \, k_1\dots k_q} - \Gamma^l_{ik_1}A^{j_1\dots j_p}_{\quad\
lk_2\dots k_q} - \dots - \Gamma^l_{ik_q}A^{j_1\dots j_p}_{\quad\
k_1\dots k_{q-1}l} .
\]
\end{definicion}

\begin{proposicionc}
Si $A^{i_1\dots i_p}_{\quad\ j_1\dots j_q}$ es un tensor de orden $(p,q)$,
entonces $\nabla_iA^{j_1\dots j_p}_{\quad\ k_1\dots k_q}$ es un tensor
de orden $(p,q+1)$.
\end{proposicionc}

\begin{proposicionc}[Compatibilidad de la métrica]
La métrica satisface
\[
\nabla_i g_{jk} = 0 .
\]
\end{proposicionc}

\section{El operador rotacional}

Trabajando en coordenadas cartesianas y usando notación de índices, el
operador rotacional se define como
\begin{align*}
\vec B = \vec\nabla\times\vec A &= E^i\partial_i\times(A^j\vec e_j) \\
&= \delta^{ik}\vec e_k \times \big[(\partial_iA^j)\vec e_j + A^j
\cancel{\partial_i\vec e_j}\big] \\
&= \vec e_k\times(\partial_iA^j)\vec e_j
= \partial_iA^j\,\vec e_i\times\vec e_j = \partial_iA^j\epsilon_{ijk}\vec
e_k ,
\end{align*}
donde el término $\partial_i\vec e_j$ se anula por trabajar en
coordenadas cartesianas ($\Gamma^k_{ij}=0$). Renombrando índices,
$\partial_jA^k\epsilon_{kij}\vec e_i = \epsilon_{ijk}\partial_jA^k\,\vec
e_i$, de donde, comparando con $\vec B = B^i\vec e_i$:

\begin{proposicionc}
En coordenadas cartesianas, el rotacional de $\vec A$ se escribe en
notación de índices como
\[
B^i = \epsilon^{ijk}\partial_jA^k .
\]
\end{proposicionc}

\section{Notación tensorial}

A partir de aquí adoptamos la \textbf{notación tensorial}: se usan los
tensores directamente, obviando los vectores base. La correspondencia es

\begin{center}
\begin{tabular}{ccc}
Notación vectorial & & Notación tensorial \\ \hline
$\vec x(x^i)$ & $\longrightarrow$ & $x^i$ \\
$\vec A(x^i)$ & $\longrightarrow$ & $A^i(x^j)$ \\
$\vec A\cdot\vec B$ & $\longrightarrow$ & $g_{ij}A^iB^j$ \\
$\vec\nabla$ & $\longrightarrow$ & $\partial/\partial x^i =: \partial_i$
\\
$\partial_iA^{j_1\dots j_p}_{\quad\ k_1\dots k_q}$ & $\longrightarrow$ &
$\nabla_iA^{j_1\dots j_p}_{\quad\ k_1\dots k_q}$
\end{tabular}
\end{center}

\begin{observacion}
En coordenadas cartesianas se recupera el caso particular
\[
\{\vec e_i\}_{i=1}^3 = \{\ihat,\jhat,\khat\}, \qquad g_{ij}=\delta_{ij},
\qquad \Gamma^k_{ij}=0, \qquad \nabla_i = \partial_i .
\]
\end{observacion}

\section{Invariancia de la acción}

Un sistema físico de $f$ grados de libertad tiene asociado un principio
de acción
\[
S = \int_{t_i}^{t_f} L\big(q^i,\dot q^i,t\big)\, dt ,
\]
del cual se deducen las ecuaciones de movimiento usando el operador
variacional,
\[
\frac{\delta L}{\delta q^i} = \frac{\partial L}{\partial q^i} -
\frac{d}{dt}\left(\frac{\partial L}{\partial \dot q^i}\right).
\]
El fundamento del principio de acción es que $S$ es un número real: se
mapea una trayectoria en el espacio de configuración a un número real.
Las ecuaciones de movimiento se obtienen exigiendo que $S$ tenga un
extremo sobre la curva que describe el movimiento físico del sistema.

Al cambiar de coordenadas, $q'^i=q'^i(q^j)$, el Lagrangiano cambiaría, en
principio, su forma funcional:
\[
S' = \int_{t_i}^{t_f} L'\big(q'^i,\dot q'^i,t\big)\, dt .
\]
Sin embargo, la curva que extremiza la acción no cambia su forma física;
solo cambian las ecuaciones (por el cambio de coordenadas) que la
describen. Por lo tanto, el valor numérico de la acción no cambia,
$S=S'$:
\[
\int_{t_i}^{t_f} L\big(q^i,\dot q^i,t\big)\, dt = \int_{t_i}^{t_f}
L'\big(q'^i,\dot q'^i,t\big)\, dt \quad\Longrightarrow\quad
L'\big(q'^i,\dot q'^i,t\big) = L\big(q^i,\dot q^i,t\big),
\]
que es la definición del nuevo Lagrangiano $L'$: el Lagrangiano es un
\textbf{escalar} bajo transformaciones de coordenadas.

\begin{ejemplo}
Consideremos el problema de fuerzas centrales conservativas, con
Lagrangiano en coordenadas cartesianas
\[
L(x,y,\dot x,\dot y) = \frac{1}{2}m(\dot x^2+\dot y^2) -
U\!\left(\sqrt{x^2+y^2}\right).
\]
Cambiando a coordenadas cilíndricas (con $z=0$), $x=r\cos\varphi$,
$y=r\sin\varphi$, se obtiene
\[
L'(r,\varphi,\dot r,\dot\varphi) = \frac{1}{2}m\left[\left(\frac{d}{dt}
(r\cos\varphi)\right)^2 + \left(\frac{d}{dt}(r\sin\varphi)\right)^2
\right] - U(r) = \frac{1}{2}m\left(\dot r^2 + r^2\dot\varphi^2\right) -
U(r).
\]
\end{ejemplo}

En una transformación de coordenadas no solo cambia la función
Lagrangiana, sino también la forma funcional de las ecuaciones de
movimiento.

\begin{proposicionc}
Bajo una transformación de coordenadas $q'^i=q'^i(q^j)$, las ecuaciones
de movimiento cambian como
\[
\frac{\delta L'}{\delta q'^i} = \frac{\partial q^j}{\partial q'^i}\,
\frac{\delta L}{\delta q^j} .
\]
Es decir, las ecuaciones de movimiento $\delta L/\delta q^i$ forman un
tensor de orden $(0,1)$.
\end{proposicionc}

